{{دیوان
| عنوان = خوشحال خاںکی وصیّت
| ماخذ = ویکی ماخذ
| ماخذ_لنک = https://ur.wikisource.org/wiki/%D8%AE%D9%88%D8%B4%D8%AD%D8%A7%D9%84_%D8%AE%D8%A7%DA%BA%DA%A9%DB%8C_%D9%88%D8%B5%DB%8C%D9%91%D8%AA
| مقابلہ = اقبال سائبر لائبریری (اقبال اکادمی پاکستان)
| مقابلہ_لنک = https://www.iqbalcyberlibrary.net
}}
<poem>
قبائل ہوں ملّت کی وحدت میں گُم
کہ ہو نام افغانیوں کا بلند

محبّت مجھے اُن جوانوں سے ہے
ستاروں پہ جو ڈالتے ہیں کمند

مغل سے کسی طرح کمتر نہیں
قہستاں کا یہ بّچۂ ارجمند

کہوں تجھ سے اے ہم نشیں دل کی بات
وہ مدفن ہے خوشحال خاں کو پسند

اُڑا کر نہ لائے جہاں بادِ کوہ
مُغل شہسواروں کی گردِ سمند!
</poem>

٭ خوشحال خاں خٹک پشتوزبان کا مشہور وطن دوست شاعر تھا جس نے افغانستان کو مغلوں سے آزادکرانے کے لیے سرحد کے افغانی قبائل کی ایک جمعیّت قائم کی۔قبائل میں صرف آفریدیوں نے آخردم تک اُس کا ساتھ دیا۔ اس کی قریبا ایک سو نظموں کا انگریزی ترجمہ ۱۸۶۲؁ء میں لندن میں شائع ہُوا تھا
